\section{Anhang}

\subsection{Styleguide}

\begin{itemize}[itemsep=1pt, parsep=0pt]
    \item \textbf{Variablen, Konstanten}
    \begin{itemize}[itemsep=1pt, parsep=0pt]
        \item mit Kleinbuchstaben beginnen
        \item erster Buchstaben von zusammengesetzten Wörtern ist gross (mixed case)
        \item keine Underscores
    \end{itemize}
    Beispiele : counter, maxSpeed
    \item \textbf{Funktionen}
    \begin{itemize}[itemsep=1pt, parsep=0pt]
        \item mit Kleinbuchstaben beginnen
        \item erster Buchstaben von zusammengesetzten Wörtern ist gross (mixed case)
        \item Namen beschreiben Tätigkeiten
        \item keine Underscores
    \end{itemize}
    Beispiele : getCount(), init(), setMaxSpeed()
    \item \textbf{Klassen, Strukturen, Enums}
    \begin{itemize}[itemsep=1pt, parsep=0pt]
        \item mit Grossbuchstaben beginnen
        \item erster Buchstaben von zusammengesetzten Wörtern ist gross (mixed case)
        \item keine Underscores
        \item Namen beschreiben Dinge
    \end{itemize}
    Beispiele : MotorController, Queue, Color
\end{itemize}

\subsection{Code beispiele}
\subsubsection{String formatting}
Utilizza gli stream di output per costruire stringhe in memoria. È sicuro perché gestisce la memoria dinamicamente.

\lstinputlisting[language=C++]{code/99_anhang/99_anhang_snippet_1.cpp}

Rappresenta lo standard moderno. Combina la sintassi concisa di \texttt{printf} con la sicurezza degli stream.

\lstinputlisting[language=C++]{code/99_anhang/99_anhang_snippet_2.cpp}

\subsubsection{Ctor un dtor}
\lstinputlisting[language=c++]{code/99_anhang/ctor_dtor_anhang.cpp}

\subsection{Ritorno per Riferimento e Method Chaining}
\begin{itemize}
    \item \textbf{Meccanismo:} Una funzione restituisce un riferimento all'oggetto stesso
    utilizzando \texttt{return *this;}. Il tipo di ritorno deve 
    essere una referenza (es. \texttt{Complex\&}).
    \item \textbf{Kaskadierung (Cascading):} Permette di chiamare più metodi in 
    sequenza sulla stessa istanza in un'unica riga: \texttt{obj.funz1().funz2();}.
    
    \crd{Tramite ritorno \textit{by-value} funz2 verrebbe eseguita su un oggetto \textit{tmp}}
    \item \textbf{Efficienza:} A differenza del ritorno per valore, 
    non viene creata alcuna copia dei dati (nessun oggetto temporaneo \textit{tmp}), 
    risparmiando risorse 
    \item \textbf{Persistenza:} Le modifiche effettuate nella catena agiscono 
    direttamente sull'oggetto originale e non su una sua copia destinata alla distruzione.
\end{itemize}

\lstinputlisting[language=c++]{code/99_anhang/method_chaining.cpp}

\subsubsection{Vergleich von Gleitkommazahlen (double)}
Beim Vergleich von \texttt{double}-Werten sollte aufgrund von Rundungsfehlern \textbf{niemals 
der Operator \texttt{==} direkt verwendet werden}. 

Stattdessen wird geprüft, ob die absolute Differenz zwischen zwei Werten kleiner als eine definierte Toleranz 
(Epsilon) ist.

\lstinputlisting[language=C++]{code/99_anhang/99_anhang_snippet_3.cpp}

\subsection{Parametri Template Non-Tipo (Static Allocation)}

I template accettano valori costanti (es. \texttt{std::size\_t N}) valutati a 
\textbf{compile-time}. Questo è l'unico modo per definire la dimensione di array statici 
(\texttt{T data[N]}) mantenendo la flessibilità del codice.

In C++, il compilatore deve conoscere l'ingombro esatto di un oggetto nello \textbf{Stack}
per generare istruzioni macchina con \textit{offset} fissi. Se \texttt{N} fosse una
variabile, il compilatore non saprebbe quanto spazio riservare automaticamente e 
l'allocazione fallirebbe.

\subsubsection{Allocazione Statica vs Dinamica}
Senza la risoluzione a tempo di compilazione, il compilatore non potrebbe riservare lo spazio nello \textbf{Stack}, obbligando all'uso della memoria \textbf{Heap} (più lenta).

\resizebox{\columnwidth}{!}{
\begin{tabular}{|l|l|l|}
\hline
\textbf{Caratteristica} & \textbf{Via Template (Statico)} & \textbf{Via Costruttore (Dinamico)} \\ \hline
\textbf{Vincolo Size}   & Costante a compile-time & Variabile a run-time \\ \hline
\textbf{Memoria}        & Stack (Veloce, fissa)   & Heap (Flessibile, lenta) \\ \hline
\textbf{Gestione}       & Automatica (No memory leaks) & Manuale (\texttt{new}/\texttt{delete}) \\ \hline
\textbf{Effetto Tipo}   & \texttt{Class<10>} $\neq$ \texttt{Class<20>} & Il tipo rimane identico \\ \hline
\end{tabular}}\vspace{1em}

\lstinputlisting[language=C++]{code/99_anhang/99_anhang_snippet_4.cpp}

\subsection{Inline Functions e Templates}

\begin{itemize}
    \item \textbf{Inlining Implicito:} Tutte le funzioni (template e non) definite direttamente all'interno della dichiarazione di una classe (\textit{in-class}) sono considerate \texttt{inline} automaticamente dal compilatore.
    \item \textbf{Inlining Esplicito:} Le funzioni template definite al di fuori della classe (\textit{out-of-class}), anche se collocate nello stesso file header, \textbf{non} sono \texttt{inline} di default. Per renderle tali, la keyword \texttt{inline} deve essere inserita tra il tag \texttt{template} e il tipo di ritorno.
    \item \textbf{Vantaggi e Limiti:} L'inlining elimina l'overhead della chiamata a funzione (\textit{zero overhead}), ma è sconsigliato per funzioni ricorsive o molto complesse.
\end{itemize}

\subsection{Overloading dell'Operatore []}

Per una corretta implementazione dell'operatore \texttt{[]}, è necessario definire \textbf{due versioni} per supportare sia gli oggetti modificabili che quelli costanti:

\begin{itemize}
    \item \textbf{Non-Const:} Restituisce una reference (\texttt{T\&}) per consentire la lettura e la scrittura.
    \item \textbf{Const:} Restituisce una const reference (\texttt{const T\&}) o un valore (\texttt{T}) per garantire il solo accesso in lettura sugli oggetti \texttt{const}.
\end{itemize}

\subsubsection{Esempio Pratico}

\lstinputlisting[language=c++]{code/99_anhang/operator.cpp}

la versione \texttt{const} viene spesso usata se la classe (l'istanza) è \texttt{const}

\subsection{Wichtige Prüfungsfallen (Trappole da Esame)}

\subsubsection{Array di Puntatori e Destruktor}
Su un array di puntatori (es. \texttt{Typ** data}), \texttt{delete[] data;} elimina solo il contenitore. Usa sempre un ciclo per deallocare i singoli oggetti ed evitare \textit{memory leaks}:
\begin{lstlisting}[language=C++]
for(int i = 0; i < groesse; ++i) { 
    delete data[i]; 
}
delete[] data;
\end{lstlisting}

\subsubsection{Default-Parameter}
I valori di default vanno definiti \textbf{esclusivamente nel file .h}. Ripeterli nel \texttt{.cpp} causa un errore di compilazione.

\subsubsection{Binding e Polimorfismo in C++}
\begin{itemize}
    \item \textbf{Chiamata diretta su oggetto} (es. \texttt{Derivata d; d.metodo();}): Nessuna ambiguità. Viene sempre risolto a compile-time (Static Binding) ed esegue il metodo della classe specifica dell'oggetto.
    \item \textbf{Chiamata tramite puntatore/reference alla classe Base} (es. \texttt{Base* p = new Derivata();}), spesso 
    succede nelle \cor{liste polimorfiche}:
    \begin{itemize}
        \item \textbf{Metodo NON \texttt{virtual} (Static/Early Binding):} La scelta del metodo dipende dal \textit{tipo del puntatore}. Verrà eseguito il metodo di \texttt{Base}.
        \item \textbf{Metodo \texttt{virtual} (Dynamic/Late Binding):} La scelta del metodo dipende dall'\textit{oggetto reale} instanziato in memoria a runtime. Verrà eseguito il metodo di \texttt{Derivata}.
    \end{itemize}
\end{itemize}

\subsection{Keyword esclusive della dichiarazione (.h)}

\resizebox{\columnwidth}{!}{
\renewcommand{\arraystretch}{1.2}
\begin{tabular}{lccc}
    \hline
    \textbf{Keyword} & \textbf{\texttt{.h}} & \textbf{\texttt{.cpp}} & \textbf{Motivo } \\
    \hline
    \texttt{static} & \checkmark & \textbf{NO} & Evita conflitti linkage\\
    \texttt{virtual} & \checkmark & \textbf{NO} & La v-table definita nell'interfaccia \\
    Parametri default & \checkmark & \textbf{NO} & Previene ambiguità \\
    \texttt{friend} & \checkmark & \textbf{NO} & Non parte della firma della funzione \\
    \texttt{explicit} & \checkmark & \textbf{NO} & direttiva di design per il compilatore \\
    \hline
\end{tabular}}

\subsection{Esempio Completo Stack}

\lstinputlisting[language=c++]{code/99_anhang/Stack.h}
\lstinputlisting[language=c++]{code/99_anhang/StackException.h}